% Section 11.4: Electron Self-Energy.
% Source: physical PDF pp. 516--520 (printed pp. 493--497).
% Starts in the lower part of p. 516 after Eq. (11.3.33); ends with Eq. (11.4.16).
% Equations: (11.4.1)--(11.4.16). Figure: 11.6. Footnotes: none.
% Status: source-reviewed against rendered pages; no unresolved uncertainties.

\subsection{Electron Self-Energy}

We conclude this chapter with a calculation of the electron self-energy
function. This by itself does not have any direct experimental implications,
but some of the results here will be useful in Chapter 14 and Volume II.

As in Section 10.3, we define $i(2\pi)^2[\Sigma^*(p)]_{\beta,\alpha}$ as the sum
of all graphs with one incoming and one outgoing electron line carrying momenta
$p$ and Dirac indices $\alpha$ and $\beta$ respectively, with the asterisk
indicating that we exclude diagrams that can be disconnected by cutting
through some internal electron line, and with propagators omitted for the two
external lines. The complete electron propagator is then given by the sum
\begin{align}
[-i(2\pi)^{-4}S'(p)]={}&[-i(2\pi)^{-4}S(p)]\notag\\
&+[-i(2\pi)^{-4}S(p)][i(2\pi)^4\Sigma^*(p)]
[-i(2\pi)^{-4}S(p)]+\cdots,
\tag{11.4.1}\label{eq:11.4.1}
\end{align}

% Figure 11.6: one-loop electron self-energy.
% Source: physical PDF p. 517 (printed p. 494).
\begingroup
\renewcommand{\thefigure}{11.\arabic{figure}}
\setcounter{figure}{5}
\begin{figure}[tbp]
\centering
\begin{tikzpicture}[
  electron/.style={line width=0.7pt,
    postaction={decorate},decoration={markings,
      mark=at position 0.52 with {\arrow{Stealth[length=4.5pt,width=3.4pt]}}}},
  photon/.style={line width=0.6pt,decorate,
    decoration={snake,amplitude=1.05pt,segment length=5pt}}
]
  \draw[electron] (-4.2,0) -- (4.2,0);
  \draw[photon] (-1.75,0) .. controls (-1.20,1.45) and (1.20,1.45) .. (1.75,0);
\end{tikzpicture}
\caption{The one-loop diagram for the electron self-energy function. As usual, the straight line represents an electron; the wavy line is a photon.}
\label{fig:11.6}
\end{figure}
\endgroup


where
\begin{equation}
S(p)=\frac{-i\sl{p}+m_e}{p^2+m_e^2-i\epsilon}.
\tag{11.4.2}\label{eq:11.4.2}
\end{equation}
The sum is trivial, and gives
\begin{equation}
S'(p)=[i\sl{p}+m_e-\Sigma^*(p)-i\epsilon]^{-1}.
\tag{11.4.3}\label{eq:11.4.3}
\end{equation}

In lowest order there is a one-loop contribution to $\Sigma^*$, given by
Figure~\ref{fig:11.6}:
\begin{align*}
i(2\pi)^4\Sigma^*_{\text{1 loop}}(p)
={}&\int d^4k\left[\frac{-i}{(2\pi)^4}
\frac{\eta_{\rho\sigma}}{k^2-i\epsilon}\right]
\cdot[(2\pi)^4e\gamma^\rho]\notag\\
&\cdot\left[\frac{-i}{(2\pi)^4}
\frac{-i\sl{p}+i\sl{k}+m_e}{(p-k)^2+m_e^2-i\epsilon}\right]
[(2\pi)^4e\gamma^\sigma]
\end{align*}
or more simply
\begin{equation}
\Sigma^*_{\text{1 loop}}(p)=\frac{ie^2}{(2\pi)^4}\int d^4k
\left[\frac{1}{k^2-i\epsilon}\right]
\left[\frac{\gamma^\rho(-i\sl{p}+i\sl{k}+m_e)\gamma_\rho}
{(p-k)^2+m_e^2-i\epsilon}\right].
\tag{11.4.4}\label{eq:11.4.4}
\end{equation}
(This is in Feynman gauge; amplitudes with charged particles off the mass
shell are not gauge-invariant.) For use in our calculation of the Lamb shift,
it will be convenient to use a method of regularization introduced by Pauli
and Villars~\hyperref[ref:11.8]{[8]}. We replace the photon propagator $(k^2-i\epsilon)^{-1}$ with
\[
\frac{1}{k^2-i\epsilon}-\frac{1}{k^2+\mu^2-i\epsilon},
\]
so that the electron self-energy function is
\begin{align}
\Sigma^*_{\text{1 loop}}(p)={}&\frac{ie^2}{(2\pi)^4}\int d^4k
\left[\frac{1}{k^2-i\epsilon}
-\frac{1}{k^2+\mu^2-i\epsilon}\right] \cdot \left[\frac{\gamma^\rho(-i\sl{p}+i\sl{k}+m_e)\gamma_\rho}
{(p-k)^2+m_e^2-i\epsilon}\right].
\tag{11.4.5}\label{eq:11.4.5}
\end{align}
Later we can drop the regulator by letting the regulator mass $\mu$ go to
infinity. In Chapter 14 we will also be interested in the case where
$\mu\ll m_e$.

We again use the Feynman trick to combine denominators, and recall that
$\gamma^\rho\gamma^\kappa\gamma_\rho=-2\gamma^\kappa$ and
$\gamma^\rho\gamma_\rho=4$. This gives
\begin{align}
\Sigma^*_{\text{1 loop}}(p)={}&\frac{ie^2}{(2\pi)^4}\int d^4k
[2i(\sl{p}-\sl{k})+4m_e]\notag\\
&\cdot\int_0^1dx\left[
\frac{1}{\{(k-px)^2+p^2x(1-x)+m_e^2x-i\epsilon\}^2}
\right.\notag\\[-2pt]
&\hspace{37mm}\left.
-\frac{1}{\{(k-px)^2+p^2x(1-x)+m_e^2x+\mu^2(1-x)-i\epsilon\}^2}
\right].
\tag{11.4.6}\label{eq:11.4.6}
\end{align}
Shifting the variable of integration $k\to k+px$ and rotating the contour of
integration gives
\begin{align}
\Sigma^*_{\text{1 loop}}(p)={}&\frac{-2\pi^2e^2}{(2\pi)^4}
\int_0^1dx\,[2i(1-x)\sl{p}+4m_e]\int_0^\infty d\kappa\,\kappa^3
\notag\\
&\cdot\left[
\frac{1}{\{\kappa^2+p^2x(1-x)+m_e^2x\}^2}
-\frac{1}{\{\kappa^2+p^2x(1-x)+m_e^2x+\mu^2(1-x)\}^2}
\right].
\tag{11.4.7}\label{eq:11.4.7}
\end{align}
The $\kappa$-integral is trivial:
\begin{equation}
\Sigma^*_{\text{1 loop}}(p)=\frac{-\pi^2e^2}{(2\pi)^4}
\int_0^1dx\,[2i(1-x)\sl{p}+4m_e]
\ln\left(\frac{p^2x(1-x)+m_e^2x+\mu^2(1-x)}
{p^2x(1-x)+m_e^2x}\right).
\tag{11.4.8}\label{eq:11.4.8}
\end{equation}

The interaction \eqref{eq:11.1.9} also contributes a renormalization counterterm
$-(Z_2-1)(i\sl{p}+m_e)+Z_2\delta m_e$ in $\Sigma^*(p)$, with $Z_2$ and
$\delta m_e$ determined by the condition that the complete propagator $S'(p)$
regarded as a function of $i\sl{p}$ should have a pole at
$i\sl{p}=-m_e$ with residue unity. (As we shall see in the next chapter,
this makes $\Sigma^*$ finite as $\mu\to\infty$ to all orders in $e$.) In
lowest order, this gives
\begin{align}
\delta m_e={}&-\left.\Sigma^*_{\text{1 loop}}\right|_{i\sl{p}=-m_e}
\notag\\
={}&\frac{2m_e\pi^2e^2}{(2\pi)^4}\int_0^1dx\,[1+x]
\ln\left(\frac{m_e^2x^2+\mu^2(1-x)}{m_e^2x^2}\right),
\tag{11.4.9}\label{eq:11.4.9}
\end{align}
\begin{align}
Z_2-1={}&-i\left.\frac{\partial\Sigma^*_{\text{1 loop}}}
{\partial\sl{p}}\right|_{i\sl{p}=-m_e}\notag\\
={}&\frac{-2\pi^2e^2}{(2\pi)^4}\int_0^1dx
\left\{(1-x)\ln\left(\frac{m_e^2x^2+\mu^2(1-x)}{m_e^2x^2}\right)
-\frac{2\mu^2(1-x)^2(1+x)}{x[m_e^2x^2+\mu^2(1-x)]}\right\}.
\tag{11.4.10}\label{eq:11.4.10}
\end{align}

(To this order, we do not distinguish between $\delta m_e$ and
$Z_2\delta m_e$.) Dropping terms that vanish for $\mu^2\to\infty$,
Eqs.~\eqref{eq:11.4.8}--\eqref{eq:11.4.10} yield
\begin{equation}
\Sigma^*_{\text{1 loop}}(p)=\frac{-\pi^2e^2}{(2\pi)^4}
\int_0^1dx\,[2i(1-x)\sl{p}+4m_e]
\ln\left(\frac{\mu^2(1-x)}{p^2x(1-x)+m_e^2x}\right),
\tag{11.4.11}\label{eq:11.4.11}
\end{equation}
\begin{equation}
\delta m_e=\frac{2m_e\pi^2e^2}{(2\pi)^4}\int_0^1dx\,[1+x]
\ln\left(\frac{\mu^2(1-x)}{m_e^2x^2}\right),
\tag{11.4.12}\label{eq:11.4.12}
\end{equation}
\begin{equation}
Z_2-1=\frac{-2\pi^2e^2}{(2\pi)^4}\int_0^1dx
\left\{(1-x)\ln\left(\frac{\mu^2(1-x)}{m_e^2x^2}\right)
-\frac{2(1-x^2)}{x}\right\}.
\tag{11.4.13}\label{eq:11.4.13}
\end{equation}
Inspection then shows that in the complete self-energy function the
$\ln\mu^2$ terms cancel, leaving us with
\begin{align}
\Sigma^*_{\text{order }e^2}(p)
={}&\Sigma^*_{\text{1 loop}}(p)-(Z_2-1)(i\sl{p}+m_e)+Z_2\delta m_e
\notag\\
={}&\frac{2\pi^2e^2}{(2\pi)^4}\int_0^1dx
\left\{[i(1-x)\sl{p}+2m_e]
\ln\left(\frac{m_e^2(1-x)}{p^2x(1-x)+m_e^2x}\right)\right.\notag\\
&\quad-m_e[1+x]\ln\left(\frac{1-x}{x^2}\right)\notag\\
&\quad\left.-(i\sl{p}+m_e)\left[(1-x)
\ln\left(\frac{1-x}{x^2}\right)-\frac{2(1-x^2)}{x}\right]\right\}.
\tag{11.4.14}\label{eq:11.4.14}
\end{align}
There is still a divergence from the behavior of the last term as $x\to0$,
which can be traced to the singular behavior of the integral over the photon
momentum $k$ in Eq.~\eqref{eq:11.4.5} at $k^2=0$, when we take $p^2$ at the point
$p^2=-m_e^2$ where we evaluated $Z_2-1$. Such infrared divergences will be
discussed in detail in Chapter 13. For the present, the point that concerns us
is that the ultraviolet divergence has cancelled.

\begin{center}
$*\ *\ *$
\end{center}

The result \eqref{eq:11.4.9} for $\delta m_e$ is of some interest in itself. Note that
$\delta m_e/m_e>0$, as we would expect for the electromagnetic self-energy due
to the interaction of a charge with its own field. But unlike the classical
estimates of electromagnetic self-energy by Poincar\'e, Abraham, and others~\hyperref[ref:11.9]{[9]},
Eq.~\eqref{eq:11.4.9} is only logarithmically divergent in the limit $\mu\to\infty$
where the cutoff is removed. In this limit:
\begin{equation}
\delta m_e\to\frac{6m_e\pi^2e^2}{(2\pi)^4}\ln\left(\frac{\mu}{m_e}\right).
\tag{11.4.15}\label{eq:11.4.15}
\end{equation}
In our calculation of the Lamb shift in Section 14.3 we will be interested in
the opposite limit, $\mu\ll m_e$. Here Eq.~\eqref{eq:11.4.9} gives
\begin{equation}
\delta m_e\to\frac{e^2\mu}{8\pi}
\left[1-\frac{\mu}{2\pi m_e}+\cdots\right].
\tag{11.4.16}\label{eq:11.4.16}
\end{equation}
